% TEMPLATE-MILC-v3.tex - plantilla maestra UNICA de las guias MILC v3.
%
% La usan las tres familias de guias, cada una con su driver:
%   · Grados 6-11  -> scripts/build-guias-g11.py
%   · Bebras       -> scripts/build-guias-bebras.py
%   · Semillero    -> scripts/build-guias-semillero.py
%
% Antes habia tres copias casi identicas (~400 lineas iguales, 6 distintas):
% cada mejora habia que aplicarla tres veces, y en la practica no se hacia
% (el motor de recursos vivia solo en dos de ellas). Lo que varia entre
% familias esta parametrizado en 6 placeholders que cada driver llena
% (se nombran aqui SIN su sintaxis real: el builder reemplaza por texto plano,
% tambien dentro de los comentarios, y un valor multilinea rompe el preambulo):
%
%   HEADER_IZQ        encabezado izquierdo de pagina
%   HEADER_DER        encabezado derecho de pagina
%   PORTADA_ETIQUETA  rotulo sobre el numero grande (GRADO/MOMENTO/MODULO)
%   PORTADA_SERIE     linea de serie de la portada
%   PORTADA_ID        identificador de la guia en la portada
%   PORTADA_CIRCULO   el circulito de grado (°), o vacio si no aplica
%
% El resto de claves las documenta content/guias/_SCHEMA.md. El script valida
% que no queden placeholders sin reemplazar antes de escribir el archivo final.

\documentclass[11pt,letterpaper]{article}
\usepackage[margin=1.75cm]{geometry}
\usepackage[table]{xcolor}
\usepackage{fontspec}
\usepackage[spanish]{babel}
\usepackage{tikz}
% Librerías TikZ para los diagramas de las guías (bloque `recursos:`):
%  · babel  → babel en español vuelve activos caracteres como «>», y TikZ los usa
%             en sus opciones (>=stealth, ->). Sin esta librería, cualquier
%             diagrama con flechas revienta con «\language@active@arg> has an
%             extra }». Debe ir ANTES de que se dibuje nada.
%  · positioning        → sintaxis «below left=2pt and 3pt».
%  · decorations.pathreplacing → llaves ({brace}) para señalar rangos.
\usetikzlibrary{babel,positioning,decorations.pathreplacing}
\usepackage{graphicx}
% Circuitos: el trabajo de electrónica se hace hoy con micro:bit y sus sensores
% integrados (MakeCode, sin cableado), así que no se necesitan esquemáticos.
% Si algún día entran montajes con protoboard, basta descomentar esta línea:
% los diagramas .tex/.tikz declarados en `recursos:` ya se incrustan solos.
% \usepackage{circuitikz}
\usepackage[most]{tcolorbox}
\usepackage{tabularx}
\usepackage{array}
\usepackage{fancyhdr}
\usepackage{titlesec}
\usepackage{ragged2e}
\usepackage{enumitem}
\usepackage{microtype}

% Helvetica Neue no trae la flecha U+2192, asi que cada «→» escrito literalmente
% en el YAML se compilaba a NADA, en silencio (462 flechas en 61 guias: rutas del
% tipo «paleta → tipografia → lockup» salian rotas). Mapeamos el caracter a su
% version matematica, que si tiene glifo. Asi el autor puede seguir escribiendo
% «→» comodamente en el YAML.
\usepackage{newunicodechar}
\newunicodechar{→}{\ensuremath{\rightarrow}}

\setmainfont{Helvetica Neue}
\setsansfont{Helvetica Neue}
\setlength{\parindent}{0pt}
\setlength{\parskip}{6pt}
\hyphenpenalty=200
\exhyphenpenalty=200
\pretolerance=400
\tolerance=2000
\setlength{\emergencystretch}{1em}
\renewcommand{\arraystretch}{1.25}
\setlist{leftmargin=5mm,itemsep=1mm,topsep=1mm}
\newcolumntype{Y}{>{\RaggedRight\arraybackslash}X}

\definecolor{milcMagenta}{HTML}{E6007E}
\definecolor{milcVino}{HTML}{5A0038}
\definecolor{milcTurquesa}{HTML}{0093A5}
\definecolor{milcVerde}{HTML}{58A923}
\definecolor{milcMostaza}{HTML}{E5B400}
\definecolor{milcOcre}{HTML}{B86600}
\definecolor{milcNegro}{HTML}{241816}
\definecolor{milcGris}{HTML}{F7F7F4}
\definecolor{milcLinea}{HTML}{E8E2DD}
\definecolor{milcGradePrimary}{HTML}{84CC16}
\definecolor{milcGradeDark}{HTML}{4D7C0F}
\definecolor{milcGradeSoft}{HTML}{ECFCCB}
\definecolor{milcGradeAccent}{HTML}{CFE7FF}
\definecolor{milcGradeText}{HTML}{FFFFFF}
\color{milcNegro}

\pagestyle{fancy}
\fancyhf{}
\lhead{\small Tecnología e Informática · Grado 7}
\rhead{\small Guía 27 · MILC v3}
\cfoot{\small\thepage}
\renewcommand{\headrulewidth}{0pt}

\newtcolorbox{titlebox}[1]{
  enhanced,colback=#1,colframe=#1,arc=7mm,boxrule=0pt,
  left=7mm,right=7mm,top=3mm,bottom=3mm,
  before skip=8pt,after skip=8pt,
  fontupper=\sffamily\bfseries\Large\color{white}
}

\newtcolorbox{softbox}[3]{
  enhanced,colback=#2,colframe=#1,borderline west={4pt}{0pt}{#1},
  arc=2mm,boxrule=.4pt,left=4mm,right=4mm,top=3mm,bottom=3mm,
  before skip=6pt,after skip=8pt,title={#3},coltitle=white,
  colbacktitle=#1,fonttitle=\sffamily\bfseries,fontupper=\RaggedRight,
  attach boxed title to top left={xshift=3mm,yshift=-2mm},
  boxed title style={arc=2mm, boxrule=0pt}
}

\newtcolorbox{drawbox}[2]{
  enhanced,colback=white,colframe=#1,arc=4mm,boxrule=1pt,
  left=5mm,right=5mm,top=4mm,bottom=4mm,before skip=8pt,after skip=8pt,
  title={#2},coltitle=white,colbacktitle=#1,fonttitle=\sffamily\bfseries,
  fontupper=\RaggedRight,
  attach boxed title to top left={xshift=4mm,yshift=-2mm},
  boxed title style={arc=3mm, boxrule=0pt}
}

\newtcolorbox{infoband}[1]{
  enhanced,colback=#1!8,colframe=#1!50,arc=2mm,boxrule=.5pt,
  left=4mm,right=4mm,top=2mm,bottom=2mm,before skip=4pt,after skip=4pt,
  fontupper=\small\RaggedRight
}

% Puente narrativo entre secciones (contrato editorial Conéctate).
% Da continuidad: "Acabas de X. Ahora vas a Y porque Z."
% Visualmente: banda izquierda en color del grado, texto en itálica pequeña.
\newtcolorbox{puentebox}{
  enhanced,colback=milcGradeSoft,colframe=milcGradeDark,
  borderline west={3pt}{0pt}{milcGradeDark},
  arc=0mm,boxrule=0pt,left=5mm,right=4mm,top=2.5mm,bottom=2.5mm,
  before skip=4pt,after skip=8pt,
  fontupper=\itshape\small\color{milcNegro}\RaggedRight
}
% La flecha va en modo matematico a proposito: Helvetica Neue NO tiene el glifo
% U+2192, asi que \textrightarrow se compilaba a NADA («Missing character: There
% is no → in font Helvetica Neue») y el puente salia sin flecha. En modo
% matematico la flecha viene de la fuente matematica y si se dibuja.
\newcommand{\puente}[1]{%
  \begin{puentebox}%
    \textcolor{milcGradeDark}{$\rightarrow$}\hspace{2mm}#1%
  \end{puentebox}%
}

\newcommand{\linead}[1]{\noindent\color{milcLinea}\rule{#1}{0.5pt}\color{milcNegro}}
\newcommand{\respuesta}[1]{\par\vspace{2mm}\foreach \n in {1,...,#1}{\noindent\color{milcLinea}\rule{\linewidth}{0.5pt}\par\vspace{5mm}}\color{milcNegro}}
\newcommand{\checkbox}{\textcolor{milcGradeDark}{\fbox{\phantom{x}}}\hspace{5pt}}

\newcommand{\phasecard}[4]{
\begin{tcolorbox}[enhanced,colback=#3,colframe=#2,arc=3mm,boxrule=.5pt,height=3.05cm,valign=center,halign=center,left=1.5mm,right=1.5mm,top=2mm,bottom=2mm]
{\sffamily\bfseries\fontsize{22}{24}\selectfont\textcolor{#2}{#1}}\par
{\sffamily\bfseries\small #4}
\end{tcolorbox}
}

% ── Recursos visuales de la guía ──────────────────────────────────────
% Declarados en el bloque `recursos:` del YAML y emitidos por el builder.
% \guiaFigura   → imagen rasterizada o PDF (foto, carta celeste, montaje).
% \guiaDiagrama → fuente TikZ/circuitikz (.tex) incrustada como vector:
%                 es la vía para circuitos y esquemáticos editables.
% #1 = ruta relativa al .tex · #2 = pie (puede ir vacío)
\newcommand{\guiaFigura}[2]{%
\begin{center}
\includegraphics[width=0.88\linewidth,height=0.38\textheight,keepaspectratio]{#1}\\[1.5mm]
{\footnotesize\itshape\textcolor{milcNegro!75}{#2}}
\end{center}
}

\newcommand{\guiaDiagrama}[2]{%
\begin{center}
\input{#1}\\[1.5mm]
{\footnotesize\itshape\textcolor{milcNegro!75}{#2}}
\end{center}
}

\begin{document}
\thispagestyle{empty}
\begin{tikzpicture}[remember picture,overlay]
  \fill[white] (current page.south west) rectangle (current page.north east);
  \path[fill=milcGradePrimary,rounded corners=16mm]
    ([xshift=.55cm,yshift=-.55cm]current page.north west) rectangle
    ([xshift=-.55cm,yshift=3.65cm]current page.south east);

  \node[anchor=north west,text=milcGradeText] at ([xshift=2.0cm,yshift=-1.85cm]current page.north west)
    {\sffamily\bfseries\fontsize{14}{16}\selectfont GRADO};
  \node[anchor=north west,text=milcGradeText,opacity=.82] at ([xshift=2.0cm,yshift=-2.75cm]current page.north west)
    {\sffamily\bfseries\fontsize{9}{11}\selectfont SERIE GUÍAS · TECNOLOGÍA E INFORMÁTICA};

  \begin{scope}[shift={([xshift=-3.05cm,yshift=-2.55cm]current page.north east)}]
    \fill[white,opacity=.80] (-.62,-.52) rectangle (.62,.52);
    \draw[milcGradeText,line width=1pt] (-.62,-.52) rectangle (.62,.52);
    \fill[milcGradeDark] (-.42,-.38) rectangle (-.22,.06);
    \fill[milcGradeAccent] (-.10,-.38) rectangle (.10,.24);
    \fill[milcGradePrimary!60!black] (.22,-.38) rectangle (.42,.38);
  \end{scope}

  \node[anchor=west,text=milcGradeText] at ([xshift=1.9cm,yshift=-12.85cm]current page.north west)
    {\sffamily\bfseries\fontsize{124}{124}\selectfont 7};
  % El circulito de grado (°) solo aplica a las guias de grado («11°»); en
  % Bebras y Semillero el numero grande es un momento/modulo, no un grado.
  \draw[milcGradeText,line width=1.7pt]
    ([xshift=4.52cm,yshift=-11.68cm]current page.north west) circle (.13cm);

  \node[anchor=west,text=milcGradeText,text width=8.7cm,align=left] at ([xshift=10.35cm,yshift=-11.6cm]current page.north west)
    {
     {\sffamily\bfseries\fontsize{19}{22}\selectfont Prompting\\avanzado ---\\técnicas que\\diferencian\\al experto}\\[4mm]
     {\sffamily\bfseries\fontsize{23}{26}\selectfont Guía 27-7-TIC}\\[2mm]
     {\sffamily\fontsize{13}{16}\selectfont Período 3 · Inteligencia Artificial}\\[5mm]
     {\sffamily\bfseries\fontsize{11}{13}\selectfont Institución Educativa Sor María Juliana}\\[1mm]
     {\sffamily\fontsize{10}{12}\selectfont Autor: Dr. Álvaro Cárdenas Orozco}
    };

  \path[fill=milcGradeSoft,rounded corners=7mm]
    ([xshift=1.35cm,yshift=.85cm]current page.south west) rectangle
    ([xshift=-1.35cm,yshift=4.25cm]current page.south east);
  \draw[milcGradeDark,line width=.65pt,rounded corners=7mm]
    ([xshift=1.35cm,yshift=.85cm]current page.south west) rectangle
    ([xshift=-1.35cm,yshift=4.25cm]current page.south east);
  \node[anchor=center,text=milcNegro,text width=17.0cm,align=center] at ([yshift=2.55cm]current page.south)
    {\sffamily\fontsize{10}{12}\selectfont Metodología MILC · Apertura ancestral · Pensamiento computacional · Triángulo Dussel-Estoicismo-Floridi\\
    \textcolor{milcGradeDark}{\bfseries Producto:} En el cuaderno: \textbf{(a) tabla de las 4 técnicas avanzadas} con definición, cuándo usarla, ejemplo; \textbf{(b) 5 prompts complejos} aplicando técnicas combinadas (mínimo 2 técnicas por prompt); \textbf{(c) prueba real} de los 5 prompts en un LLM con la respuesta obtenida; \textbf{(d) tu evaluación} de cuál técnica fue más útil. Firma + reflexión personal.};
\end{tikzpicture}
\mbox{}
\newpage

\section*{Apertura: Prompting avanzado --- técnicas que diferencian al usuario experto}

\begin{softbox}{milcGradeDark}{milcGradeSoft}{Saber ancestral}
\textbf{Don Lucho el relojero del Valle del Cauca tenía un truco para enseñar a sus aprendices a reparar relojes.} Cuando uno de ellos --- digamos un muchacho de 14 años de la vereda --- llegaba con un reloj averiado, don Lucho NO le decía: \emph{``arréglalo''}. Le decía: \emph{``imagínate que tú eres yo. ¿Qué harías primero? ¿Por qué? Explícame paso a paso''}. El aprendiz pensaba en voz alta y don Lucho lo corregía suavemente. \textbf{Después de 2 años, el aprendiz hacía esa misma reflexión solo}: se hablaba a sí mismo como si fuera don Lucho. Esa técnica antigua --- \emph{adoptar el rol de un experto, pensar en voz alta, mostrar tus pasos} --- es exactamente lo que en prompting moderno se llama \textbf{chain-of-thought} y \textbf{role prompting}. Las técnicas avanzadas no son inventos de Silicon Valley: son sabidurías pedagógicas antiguas \emph{trasladadas a conversaciones con IA}.
\end{softbox}

\begin{softbox}{milcTurquesa}{milcGris}{Saber contemporáneo}
\textbf{El prompting avanzado} usa técnicas que afinan las respuestas del LLM mucho más allá de la fórmula CTRF básica. Las \textbf{4 técnicas más útiles para 7° grado:} \textbf{(1) Role Prompting} (asignar un rol): \emph{``Actúa como un profesor de matemáticas paciente que enseña a niños de 12 años''}. \textbf{(2) Few-shot} (dar ejemplos): \emph{``Te muestro 2 ejemplos de cómo quiero que respondas, después tú haces el 3°''}. \textbf{(3) Chain-of-Thought} (pensar en voz alta): \emph{``Resuelve este problema mostrando todos los pasos de tu razonamiento, no solo la respuesta final''}. \textbf{(4) Iteración explícita} (pedir mejoras): \emph{``Primero hazlo así. Después dime cómo mejorarlo. Después implementa la mejora''}. \textbf{Estas técnicas se pueden combinar}: \emph{Role + Chain-of-Thought = pídele que actúe como experto Y muestre su razonamiento}. La combinación es donde está la verdadera magia del prompting avanzado.
\end{softbox}

\begin{softbox}{milcMostaza}{milcGris}{Pregunta-puente}
\textit{Si le pides a ChatGPT \emph{``actúa como un profesor estricto de matemáticas y resuelve este problema mostrando todos los pasos''}, ¿\textbf{cómo crees que cambia} la respuesta vs simplemente preguntar el problema?}
\end{softbox}

\begin{softbox}{milcVerde}{milcGris}{Saber y saber hacer}
Al terminar la clase: (1) podrás \textbf{identificar} las 4 técnicas avanzadas; (2) sabrás \textbf{aplicar} cada una a una situación; (3) podrás \textbf{combinar} 2 o más técnicas en un mismo prompt; (4) habrás \textbf{creado} 5 prompts complejos y probados con LLM.
\end{softbox}



\small
\begin{tabularx}{\textwidth}{>{\bfseries}p{3.0cm}Y}
\rowcolor{milcGradePrimary!12}Autor & Dr. Álvaro Cárdenas Orozco\\
DBA & Reconoce los fundamentos, usos y límites éticos de la Inteligencia Artificial (MEN, T\&I 7°)\\
Referentes & Saber ancestral del consejero · Modelos de lenguaje · Ética algorítmica y prompting\\
Producto final & En el cuaderno: \textbf{(a) tabla de las 4 técnicas avanzadas} con definición, cuándo usarla, ejemplo; \textbf{(b) 5 prompts complejos} aplicando técnicas combinadas (mínimo 2 técnicas por prompt); \textbf{(c) prueba real} de los 5 prompts en un LLM con la respuesta obtenida; \textbf{(d) tu evaluación} de cuál técnica fue más útil. Firma + reflexión personal.\\
\end{tabularx}
\normalsize

\puente{Hoy haces 4 cosas: descubres 4 técnicas avanzadas, las aplicas a ejemplos, las combinas, las pruebas.}

\begin{titlebox}{milcGradeDark}
Ruta MILC: así vamos a aprender
\end{titlebox}
\begin{tabularx}{\textwidth}{>{\centering\arraybackslash}X>{\centering\arraybackslash}X>{\centering\arraybackslash}X>{\centering\arraybackslash}X}
\phasecard{1}{milcVerde}{milcGris}{Escucha\\[-1mm]\small Observo, escucho y pregunto.} &
\phasecard{2}{milcTurquesa}{milcGris}{Sistematización\\[-1mm]\small Pensamiento computacional.} &
\phasecard{3}{milcMagenta}{milcGris}{Praxis\\[-1mm]\small Construyo y aplico.} &
\phasecard{4}{milcVino}{milcGris}{Evaluación\\[-1mm]\small Reflexión liberadora.}
\end{tabularx}


\puente{Empezamos viendo cómo cambia la respuesta del LLM con técnica vs sin técnica.}

\begin{titlebox}{milcVerde}
Fase 1 · Escucha
\end{titlebox}

\begin{softbox}{milcVerde}{milcGris}{Escena para escuchar}
\textbf{Actividad 1 · ANALIZA --- Compara 2 versiones de un prompt} (10 min · individual). Lee estos 2 prompts y en el cuaderno responde \emph{¿cuál crees que dará mejor respuesta? ¿Por qué?}. \textbf{Versión A (CTRF básico):} \emph{``Soy estudiante de 7°. Resuelve este problema: si Pedro tiene 24 manzanas y reparte la cuarta parte a 4 amigos, ¿cuántas le quedan? Hazlo paso a paso''}. \textbf{Versión B (con técnicas avanzadas):} \emph{``Actúa como un profesor paciente de matemáticas de secundaria (role). Resuelve este problema mostrando TODOS los pasos de tu razonamiento, como si lo estuvieras enseñando a un niño que no entiende muy bien (chain-of-thought): si Pedro tiene 24 manzanas y reparte la cuarta parte a 4 amigos, ¿cuántas le quedan? Después de resolverlo, explícame también qué propiedad matemática usé y dame un ejercicio parecido para practicar (iteración explícita)''}.
\end{softbox}

\checkbox Leí las 2 versiones con calma.\\
\checkbox Identifiqué cuál tiene técnicas avanzadas.\\
\checkbox Predigo que B dará mejor respuesta porque es más completa.

\begin{infoband}{milcVerde}


\textbf{Lo que va al cuaderno (Actividad 1 · ANALIZA).} Título: «Actividad 1 --- A vs B: con y sin técnicas». \textbf{Formato:} respuesta corta sobre cuál mejor y por qué. \textbf{Extensión:} 4 líneas. \textbf{Sabes que terminaste cuando} reconoces que B tiene rol + chain-of-thought + iteración explícita.
\end{infoband}

\puente{Después aprendes las 4 técnicas con ejemplos detallados.}

\begin{titlebox}{milcTurquesa}
Fase 2 · Sistematización con pensamiento computacional
\end{titlebox}

\begin{softbox}{milcTurquesa}{milcGris}{Cuatro pilares aplicados al tema}
La versión B sería mejor porque combina 3 técnicas. Ahora aprende las \textbf{4 técnicas avanzadas} a fondo y cuándo usar cada una.
\end{softbox}

\small
\begin{tabularx}{\textwidth}{>{\bfseries\RaggedRight\arraybackslash}p{3.6cm}Y}
\rowcolor{milcTurquesa!90}\color{white}Pilar & \color{white}\bfseries Aplicación al tema\\
1. Descomposición & \textbf{Técnica 1 --- Role Prompting (asignar un rol).} Le dices a la IA que actúe como X experto. \emph{Ejemplos:} \emph{``Actúa como un médico explicando a un paciente''}, \emph{``Actúa como un poeta del siglo XIX''}, \emph{``Actúa como un programador de Python''}. \textbf{¿Por qué funciona?} El LLM ajusta tono, vocabulario y profundidad según el rol. \emph{Útil para}: explicaciones técnicas, escritura creativa, simulaciones. \emph{Cuidado}: no le pidas que actúe como persona real (puede inventar datos atribuidos a esa persona).\\
2. Reconocimiento de patrones & \textbf{Técnica 2 --- Few-shot (dar ejemplos).} Le muestras 1-3 ejemplos de cómo quieres que responda, después le pides que haga uno nuevo siguiendo el patrón. \emph{Ejemplo}: \emph{``Te muestro 2 ejemplos de cómo quiero que clasifiques canciones. Ejemplo 1: 'La gota fría' = vallenato. Ejemplo 2: 'Despacito' = reggaetón. Ahora clasifica: 'La vida es un carnaval'''}. \textbf{¿Por qué funciona?} El LLM aprende el patrón de los ejemplos y lo aplica. \emph{Útil para}: tareas con formato específico, clasificaciones, traducciones especiales.\\
3. Abstracción & \textbf{Técnica 3 --- Chain-of-Thought (pensar en voz alta).} Le pides que muestre el razonamiento paso a paso, no solo la respuesta final. \emph{Ejemplo}: \emph{``Resuelve este problema mostrando TODOS los pasos de tu razonamiento, como si me estuvieras enseñando''}. \textbf{¿Por qué funciona?} El LLM se ``obliga'' a pensar antes de responder. Las respuestas son más precisas (especialmente en matemáticas, lógica, análisis). \emph{Útil para}: problemas de matemáticas, ejercicios de lógica, análisis complejos.\\
4. Algoritmo conceptual & \textbf{Técnica 4 --- Iteración explícita (pedir mejoras).} En lugar de aceptar la primera respuesta, le pides al LLM que se autocritique y mejore. \emph{Ejemplo}: \emph{``Primero respóndeme. Después dime cuáles son los 3 puntos débiles de tu respuesta. Después dame una versión mejorada que aborde esos puntos''}. \textbf{¿Por qué funciona?} El LLM hace ``autoevaluación'' y produce versión más sólida. Genera respuestas que normalmente requerirían 3-4 mensajes en 1 solo. \emph{Útil para}: ensayos, decisiones importantes, planes complejos. \textbf{Combinaciones poderosas:} \emph{Role + Chain-of-Thought}: experto pensando en voz alta. \emph{Few-shot + Iteración}: ejemplos + autocrítica. \emph{Role + Few-shot}: experto siguiendo formato específico.\\
\end{tabularx}
\normalsize

\begin{drawbox}{milcTurquesa}{4 técnicas avanzadas + combinaciones}
\textbf{4 técnicas:} \\[1mm]
\textbf{1.} ROLE: actúa como [experto]. \\[1mm]
\textbf{2.} FEW-SHOT: te doy ejemplos, sigue el patrón. \\[1mm]
\textbf{3.} CHAIN-OF-THOUGHT: muestra todos los pasos. \\[1mm]
\textbf{4.} ITERACIÓN EXPLÍCITA: dame respuesta, critícala, mejórala. \\[1mm]
\textbf{Combinaciones poderosas:} Role + CoT · Few-shot + Iteración · Role + Few-shot.
\end{drawbox}

\begin{softbox}{milcMostaza}{milcGris}{Errores comunes y su corrección}
\textbf{Errores frecuentes al usar técnicas avanzadas.} (1) \emph{Usar demasiadas técnicas a la vez:} prompt de 500 palabras es contraproducente. Combina 2-3 máximo. (2) \emph{Role inadecuado:} pedir que actúe como una persona real, como un dios, como algo imposible. La IA puede inventar. (3) \emph{Few-shot con malos ejemplos:} si los ejemplos son malos, el resultado también. (4) \emph{Chain-of-Thought esperando que la IA tenga razón siempre:} muestra pasos, pero verifica las matemáticas. (5) \emph{Iteración sin paciencia:} si pides ``mejóralo'' 10 veces seguidas, las mejoras se vuelven cosméticas.
\end{softbox}

\begin{infoband}{milcTurquesa}


\textbf{Lo que va al cuaderno (Actividad 2 · EVALÚA).} Título: «Actividad 2 --- Qué técnica de prompting es la mejor para cada tarea». \textbf{Formato:} tabla de 4 filas, 3 columnas (técnica, cuándo usarla, ejemplo). + 3 combinaciones poderosas. \textbf{Veredicto (3 tareas):} para estas 3 tareas decide qué técnica (o combinación) usarías y justifica con 1 razón: (1) resolver un problema de física paso a paso; (2) imitar el estilo de redacción de tu profe favorito; (3) generar 5 variantes de un correo sin sonar a IA. \textbf{Extensión:} 7 líneas + veredicto. \textbf{Sabes que terminaste cuando} puedes recomendar la técnica correcta para cualquier tarea sin titubear.
\end{infoband}

\puente{Con todo claro, creas 5 prompts complejos y los pruebas.}

\begin{titlebox}{milcMagenta}
Fase 3 · Praxis con pensamiento computacional
\end{titlebox}

\begin{softbox}{milcMagenta}{milcGris}{Construyendo el producto con los cuatro pilares}
\textbf{Actividad 3 · APLICA --- 5 prompts complejos con técnicas combinadas} (32 min · individual). Vas a crear 5 prompts complejos aplicando 2+ técnicas y probarlos con un LLM real.
\end{softbox}

\small
\begin{tabularx}{\textwidth}{>{\bfseries\RaggedRight\arraybackslash}p{3.6cm}Y}
\rowcolor{milcMagenta!90}\color{white}Pilar & \color{white}\bfseries Aplicación al producto\\
Descomposición & \textbf{Paso 1 --- Plan de los 5 prompts.} En el cuaderno escribe 5 temas (cosas reales que quieras consultar). Sugerencias: \emph{(1) Resolver un problema de mate con explicación}; \emph{(2) Pedir consejo sobre cómo estudiar para examen}; \emph{(3) Recibir crítica constructiva sobre algo que escribiste}; \emph{(4) Aprender un concepto científico desde 2 ángulos}; \emph{(5) Planear una sorpresa para alguien}.\\
Patrones & \textbf{Paso 2 --- Diseña los 5 prompts combinando técnicas.} Para cada uno, escribe el prompt completo con CTRF + 2 técnicas avanzadas. \emph{Ejemplo para Prompt 1}: \emph{``Soy estudiante de 7° en Colombia (C). Actúa como un profesor paciente de matemáticas (Role) y resuelve este problema mostrando TODOS los pasos del razonamiento (Chain-of-Thought): [problema]. Después del razonamiento, dame 2 errores típicos que los estudiantes cometen al resolver problemas así (R + F)''}. Aplica creatividad.\\
Abstracción & \textbf{Paso 3 --- Tabla en el cuaderno.} Tabla de \textbf{5 filas y 4 columnas}: \emph{(1) Tema del prompt}, \emph{(2) Prompt completo}, \emph{(3) Técnicas usadas}, \emph{(4) Resultado obtenido (resumen)}. Llena las columnas 1, 2 y 3.\\
Algoritmo de armado & \textbf{Paso 4 --- Prueba los 5 prompts en un LLM.} Abre el LLM (Claude, ChatGPT, Gemini). Copia cada prompt completo. Lee la respuesta. Anota en columna 4 un resumen de lo que el LLM produjo.\\
\end{tabularx}
\normalsize

\begin{drawbox}{milcMagenta}{Antes de cerrar la S7}
\checkbox La tabla tiene 5 prompts complejos completos. \\[1mm]
\checkbox Cada prompt aplica mínimo 2 técnicas avanzadas. \\[1mm]
\checkbox Probé los 5 prompts en un LLM real. \\[1mm]
\checkbox Anoté resumen de las 5 respuestas. \\[1mm]
\checkbox Hay evaluación de cuál técnica fue más útil. \\[1mm]
\checkbox La reflexión final está firmada.
\end{drawbox}

\begin{softbox}{milcMagenta}{milcGris}{Plantilla de trabajo}
\textbf{Estructura.} \textit{Paso 1:} 5 temas. \textit{Paso 2:} diseñar 5 prompts con técnicas combinadas. \textit{Paso 3:} tabla 5x4. \textit{Paso 4:} probar en LLM. \textit{Paso 5:} evaluar técnicas. \textit{Paso 6:} reflexión + firma.
\end{softbox}

\begin{infoband}{milcMagenta}


\textbf{Lo que va al cuaderno (Actividad 3 · APLICA).} Título: «Actividad 3 --- 5 prompts complejos con técnicas combinadas». \textbf{Formato:} tabla 5x4 + evaluación + reflexión. \textbf{Extensión:} 1 página y media. \textbf{Sabes que terminaste cuando} eres usuario avanzado de prompting.
\end{infoband}

\puente{El producto son los 5 prompts probados + tabla + evaluación.}

\begin{titlebox}{milcMostaza}
Producto de la sesión
\end{titlebox}
\textbf{5 prompts complejos con técnicas combinadas · probados con LLM · evaluados · firmado}.
\begin{softbox}{milcMostaza}{milcGris}{Criterios de calidad}
(1) Hay 5 prompts complejos completos. (2) Cada prompt aplica mínimo 2 técnicas avanzadas. (3) Los 5 fueron probados con un LLM real. (4) Hay evaluación de cuál técnica fue más útil. (5) Hay reflexión final personal firmada.
\end{softbox}

\puente{Antes de salir, evalúa cuál técnica funcionó mejor para qué tipo de tarea.}

\begin{titlebox}{milcVino}
Fase 4 · Evaluación liberadora — cinco dimensiones
\end{titlebox}

\begin{softbox}{milcVino}{milcGris}{Sentido de la evaluación}
Evaluar no es solo revisar una nota. Es reconocer qué aprendí, qué cambió en mí, qué emociones manejé, cómo me relaciono con la ciudad y la comunidad, y qué le dejaré a quien viene detrás.
\end{softbox}

\textbf{Mi reflexión liberadora en cinco dimensiones.} Completa cada frase iniciada:

\small
\begin{tabularx}{\textwidth}{>{\bfseries\RaggedRight\arraybackslash}p{3.4cm}Y>{\RaggedRight\arraybackslash}p{4.6cm}}
\rowcolor{milcVino}\color{white}Dimensión & \color{white}\bfseries Mi respuesta (frame starter) & \color{white}\bfseries Algo que descubrí\\
1. Personal & Lo que más cambió en mí fue \linead{6.5cm} & \linead{4.2cm}\\
2. Emocional & La emoción más fuerte fue \linead{2.5cm} y la regulé \linead{3cm} & \linead{4.2cm}\\
3. Ciudadana & En democracia esto importa porque \linead{6cm} & \linead{4.2cm}\\
4. Local (Cartago) & En mi barrio sería útil para \linead{6.5cm} & \linead{4.2cm}\\
5. Intergeneracional & A mi abuela/abuelo le diría \linead{6.5cm} & \linead{4.2cm}\\
\end{tabularx}
\normalsize

\begin{infoband}{milcVino}
\textbf{Para la próxima sustentación.} Vuelve a esta tabla en seis meses. Verás que las respuestas cambian — no porque lo aprendido cambie, sino porque tú cambias. La evaluación liberadora es una conversación contigo a través del tiempo.
\end{infoband}


\puente{Cerramos con 3 voces sobre por qué el prompting es habilidad nueva clave.}

\begin{titlebox}{milcGradeDark}
Triángulo de pensamiento · cierre filosófico
\end{titlebox}

\begin{softbox}{milcVino}{milcGris}{Dussel · lente del nosotros}
\textit{````El que pasa de aficionado a artesano lo hace combinando técnicas. La calidad nace de la combinación, no de la adición.''''}

Don Lucho el relojero no era bueno solo por usar herramientas: era bueno por \textbf{combinarlas con criterio}. Lupa + paciencia + delicadeza + experiencia. \emph{Lo mismo con prompting: las técnicas individuales son herramientas; la combinación con criterio es oficio}. Hoy diste el paso de aficionado (CTRF básico) a aprendiz de artesano (combinación de técnicas). El oficio sigue por años más, pero ya entraste al taller.

\textbf{¿Estoy combinando las herramientas con criterio, o usándolas una por una sin pensar?}
\end{softbox}
\linead{\linewidth}\\[1mm]
\linead{\linewidth}

\begin{softbox}{milcMostaza}{milcGris}{Estoicismo · Marco Aurelio · cuidado interior}
\textit{````El sabio no se conforma con saber lo básico. Profundiza hasta dominar las sutilezas. Ahí está la diferencia con los demás.''''}

Mucha gente aprende solo CTRF básico (la S6) y se conforma. \textbf{Tú aprendes técnicas avanzadas (S7)}: te coloca en grupo más pequeño de personas que dominan prompting. \emph{Esa profundización en lo que aprendes es disciplina antigua}: \emph{no satisfacerse con lo básico cuando hay más por aprender}. La sutileza es donde está la maestría.

\textbf{¿En qué áreas de mi vida me conformo con lo básico cuando podría profundizar?}
\end{softbox}
\linead{\linewidth}\\[1mm]
\linead{\linewidth}

\begin{softbox}{milcTurquesa}{milcGris}{Floridi · lente de la infoesfera}
\textit{````Las habilidades que valen oro en 2030 son las que pocos dominan en 2026. El prompting avanzado es una de ellas.''''}

\textbf{En 2030, las personas que dominen prompting avanzado tendrán ventaja en universidad y trabajo}. Hoy son pocos los que saben combinar Role + Chain-of-Thought con criterio. Tú lo aprendiste en grado 7. \emph{Esa diferencia se acumula con los años}. Mientras tus pares descubren CTRF básico a los 20, tú ya estarás en niveles más avanzados.

\textbf{¿Estoy aprovechando este momento clave para construir habilidades que pocos tendrán?}
\end{softbox}
\linead{\linewidth}\\[1mm]
\linead{\linewidth}

\begin{titlebox}{milcVino}
Compromiso final
\end{titlebox}

\small
\begin{tabularx}{\textwidth}{>{\RaggedRight\arraybackslash}p{3.0cm}YYY}
\rowcolor{milcVino}\color{white}\bfseries Criterio & \color{white}\bfseries Lo logré & \color{white}\bfseries En proceso & \color{white}\bfseries Necesito apoyo\\
Comprensión del tema & Explico ideas centrales con mis palabras. & Reconozco ideas pero falta precisión. & Necesito repasar conceptos base.\\
Pensamiento computacional & Apliqué los 4 pilares al diseñar mi producto. & Apliqué algunos pilares. & Necesito releer la fase 2.\\
Aplicación y producto & Mi producto cumple los criterios y se puede revisar. & Producto funcional con ajustes pendientes. & Aún en construcción.\\
Reflexión 5 dimensiones & Respondí cada dimensión con honestidad. & Respondí algunas con detalle. & Mis respuestas son muy generales.\\
Triángulo de pensamiento & Conecté las 3 voces con mi trabajo real. & Conecté al menos una. & No relaciono las voces con mi proceso.\\
\end{tabularx}
\normalsize

\textbf{Hoy entendí que:}\\[1mm]
\linead{\linewidth}\\[1mm]
\linead{\linewidth}

\textbf{Mi compromiso para la próxima sesión es:}\\[1mm]
\linead{\linewidth}\\[1mm]
\linead{\linewidth}

\begin{softbox}{milcMostaza}{milcGris}{Cita de despedida · Marco Aurelio}
\textit{``El obstáculo en el camino se vuelve el camino. Lo que estorba al camino se vuelve camino.''}\\[1mm]
Cada error de hoy, bien observado, se convierte en pieza del camino que pisarás mañana.
\end{softbox}

\small
\begin{tabularx}{\textwidth}{>{\bfseries}p{3.6cm}Y}
\rowcolor{milcGradeDark!12}Nombre completo: & \linead{10cm}\\
Fecha: & \linead{4cm} \quad Curso/grupo: \linead{4cm}\\
Mi firma: & \linead{9cm}\\
\end{tabularx}
\normalsize

\begin{infoband}{milcGradeDark}
\textbf{Para la próxima sesión.} Esta semana, usa al menos 1 técnica avanzada en cada conversación con LLM. La próxima clase aprendes \textbf{ética de la IA}: sesgo, privacidad, derechos de autor, deepfakes. Es la otra cara fundamental del uso responsable.
\end{infoband}

\end{document}
